%% Работа с русским языком
\usepackage{cmap}			 % поиск в PDF
% mathtext убран: в формулах нет «голой» кириллицы (вся русская — в \text{...}),
% а с tempora (6 кодировок fontenc) mathtext переполнял лимит 16 math-alphabet.
\usepackage[T2A]{fontenc}	 % кодировка
\usepackage[utf8]{inputenc}	 % кодировка исходного текста
\usepackage[english,russian]{babel}	 % локализация и переносы

%% Пакеты для работы с математикой
\usepackage{amsmath,amsfonts,amssymb,amsthm,mathtools}
\usepackage{icomma}

%% Нумерация формул (опционально)
%\mathtoolsset{showonlyrefs=true} % показывать номера только у тех формул, на которые есть \eqref{} в тексте.
%\usepackage{leqno}               % нумерация формул слева

%% Шрифты
\usepackage{tempora}     % Times-подобный шрифт с кириллицей (ГОСТ) + matching math

%% Некоторые полезные макросы для дебага (в случае недоверия авторам шаблона)
\makeatletter
\newcommand\thefontsize{The current font size is: \f@size pt} % пример: \section{\thefontsize}
\makeatother

%% Настройка размеров шрифтов
%% Тело — 14 pt (класс extarticle). Заголовки тоже 14 pt: по ГОСТ они того же
%% кегля, что и текст, и выделяются полужирным начертанием.
\makeatletter
\setlength{\headheight}{28pt}
\renewcommand\Huge{\@setfontsize\Huge{14pt}{17pt}}
\renewcommand\huge{\@setfontsize\huge{14pt}{17pt}}
\renewcommand\Large{\@setfontsize\Large{14pt}{17pt}}
\renewcommand\large{\@setfontsize\large{14pt}{17pt}}
\makeatother

%% Поля (геометрия страницы)
\usepackage[left=3cm,right=1.5cm,top=2cm,bottom=2cm,bindingoffset=0cm]{geometry}

%% Русские списки
\usepackage{enumitem}
\makeatletter
\AddEnumerateCounter{\asbuk}{\russian@alph}{щ}
\makeatother

%% Работа с картинками
\usepackage{caption} % пакет для подписей к рисункам, в частности, для работы caption*
\captionsetup{justification=centering} % центрирование подписей к картинкам
\usepackage{graphicx}                  % вставки рисунков
\graphicspath{{images/}{images2/}}     % папки с картинками
\setlength\fboxsep{3pt}                % отступ рамки \fbox{} от рисунка
\setlength\fboxrule{1pt}               % толщина линий рамки \fbox{}
\usepackage{wrapfig}                   % обтекание рисунков и таблиц текстом
\usepackage{placeins}                  % \FloatBarrier — управление размещением плавающих объектов

%% Работа с таблицами
\usepackage{array,tabularx,tabulary,booktabs} % дополнительная работа с таблицами
\usepackage{longtable}                        % длинные таблицы
\usepackage{multirow}                         % слияние строк и столбцов в таблице
\usepackage{makecell}                         % создание ячеек с многими строчками текста
\usepackage{hhline}                           % позволяет разрывать горизонтальные линии в таблицах

%% Красная строка
\setlength{\parindent}{1.25cm}

%% Интервалы (полуторный межстрочный по ГОСТ)
\usepackage{setspace}
\onehalfspacing

%% TikZ (структурные схемы узлов)
\usepackage{tikz}
\usetikzlibrary{graphs,graphs.standard,arrows.meta,calc,positioning}

%% Колонтитулы: по ГОСТ только номер страницы внизу по центру,
%% без верхнего колонтитула и без линеек
\usepackage{fancyhdr}
\pagestyle{fancy}
\fancyhf{}
\fancyfoot[C]{\thepage}
\renewcommand{\headrulewidth}{0pt}
\renewcommand{\footrulewidth}{0pt}

%% Перенос знаков в формулах (по Львовскому)
\newcommand*{\hm}[1]{#1\nobreak\discretionary{}{\hbox{$\mathsurround=0pt #1$}}{}}

%% Дополнительно
\usepackage{float}   % добавляет возможность работы с командой [H] которая улучшает расположение на странице
\usepackage{gensymb} % красивые градусы
\usepackage{listings} % пакет для листингов с кодом
\lstset{              % настройки для лисингов с кодом
basicstyle=\small\ttfamily,
columns=flexible,
breaklines=true
}

% Микротипографика (улучшенный перенос и трекинг)
\usepackage{microtype}

% Допуск на растяжение межсловных пробелов в строках с неразрывными длинными
% токенами (артикулы вида SN74HC4066N, шестнадцатеричные коды) — устраняет вылет за поле.
\emergencystretch=3em

% Вставка внешнего PDF-титула
\usepackage{pdfpages}

% Hyperref (для ссылок внутри  pdf)
\usepackage[unicode, pdftex, hidelinks]{hyperref}

% Отступ перед первым абзацем в каждом разделе
\usepackage{indentfirst}
